\begin{center}
    {\LARGE \textbf{2009物理化学}} \\[2ex]
    {\large 时量180分钟 \quad 满分90分}
\end{center}

\vspace{0.5cm}
\noindent\textbf{注意事项:}
\begin{enumerate}[label=\arabic*., parsep=0pt, itemsep=0pt]
    \item 答题前,考生先将自己的姓名、学号、考场号、座位号填写在答题卡上,将条形码准确粘贴在考生信息条形码粘贴区。
    \item 考试结束后,将本试卷与答题纸一并收回。
    \item 回答各个问题时,将答案写在答题纸上,写在本试卷上无效。
\end{enumerate}

\vspace{0.5cm}
\noindent\textbf{1. (15 pts.)}

取$0^\circ\mathrm{C}$,$3p^\ominus$的$\ce{O_{2}(g)}$ $10\ \mathrm{dm^3}$,绝热膨胀到压力为$p^\ominus$,分别计算下列两种过程的$\Delta G$:
\begin{enumerate}[label=(\arabic*), itemsep=1ex]
    \item 绝热可逆膨胀；
    \item 将外压力骤减至$p^\ominus$,气体反抗外压力进行绝热不可逆膨胀。
\end{enumerate}
假定$\ce{O_{2}(g)}$为理想气体,其摩尔定容热容$C_{V,\mathrm{m}}=\dfrac{5}{2}R$。已知氧气的摩尔标准熵$S_\mathrm{m}^\ominus(298\ \mathrm{K})=205.0\ \mathrm{J\cdot K^{-1}\cdot mol^{-1}}$。

\vspace{0.5cm}
\noindent\textbf{2. (15 pts.)}

人类血浆的凝固点为$272.66\ \mathrm{K}$,求$310.2\ \mathrm{K}$时血浆的渗透压。已知水的凝固点降低常数$K_\mathrm{f}=1.86\ \mathrm{K\cdot mol^{-1}\cdot kg}$。

\vspace{0.5cm}
\noindent\textbf{3. (15 pts.)}

苯乙烯在高温下易聚合,故苯乙烯的精制采用减压蒸馏,若控制蒸馏温度为$303\ \mathrm{K}$,试求蒸馏设备需要达到的真空度。

已知苯乙烯的沸点$T_\mathrm{b}=418\ \mathrm{K}$,蒸发热$\Delta_\mathrm{vap}H_\mathrm{m}=40.313\ \mathrm{kJ\cdot mol^{-1}}$。

\vspace{0.5cm}
\noindent\textbf{4. (15 pts.)}

液态$\ce{Br_{2}}$在$331.4\ \mathrm{K}$时沸腾,$\ce{Br_{2}(l)}$在$282.5\ \mathrm{K}$时的蒸气压为$13.33\ \mathrm{kPa}$,

计算$298.15\ \mathrm{K}$时$\ce{Br_{2}(g)}$的标准生成吉布斯自由能。

\vspace{0.5cm}
\noindent\textbf{5. (15 pts.)}

在$433\ \mathrm{K}$时,气相反应$\ce{N_{2}O_{5} -> 2NO_{2} + \dfrac{1}{2}O_{2}}$是一级反应
\begin{enumerate}[label=(\arabic*), itemsep=1ex]
    \item 在等容容器中最初引入纯的$\ce{N_{2}O_{5}}$,$3\ \mathrm{s}$后容器压力增大一倍。
    \begin{enumerate}[label=\textcircled{\arabic*}, itemsep=1ex]
        \item 求此时$\ce{N_{2}O_{5}}$的分解百分数；
        \item 求速率常数。
    \end{enumerate}
    \item 若反应发生在同样容器中但温度为$T_2$,在$3\ \mathrm{s}$后容器的压力增大到最初的$1.5$倍。
    \begin{enumerate}[label=\textcircled{\arabic*}, itemsep=1ex]
        \item 求温度$T_2$时反应的半衰期；
        \item 求温度$T_2$(已知活化能是$103\ \mathrm{kJ\cdot mol^{-1}}$)。
    \end{enumerate}
\end{enumerate}

\vspace{0.5cm}
\noindent\textbf{6. (15 pts.)}

$25^\circ\mathrm{C}$时,电池$\ce{Pt|H_{2}(p^\ominus)|H_{2}SO_{4}(4.00\ mol\cdot kg^{-1})|Hg_{2}SO_{4}(s)|Hg(l)}$的电动势$E=0.6120\ \mathrm{V}$,

其标准电动势$E^\ominus=0.6152\ \mathrm{V}$,求该$\ce{H_{2}SO_{4}}$溶液的离子平均活度因子$\gamma_\pm$。